% !TEX root = ../../main.tex

\subsection{Edge Multiplicity Control Across the Tree}

Even a well-calibrated edge test would over-split the hierarchy if it were
applied independently to every edge and thresholded at a fixed level. The tree
contains many overlapping candidate splits, and the final partition depends on
reading those edge decisions jointly. The method therefore applies a
hierarchical Benjamini--Hochberg correction on the tree
\citep{bogomolov2021hypotheses} at level
$\alpha_{\mathrm{edge}} = 0.001$. Edge hypotheses are grouped into families by
parent node. If $F(u)$ denotes the family of edges from parent $u$ to its
children, then the family-specific level used by the implementation is
\begin{equation}
\alpha_u
=
\alpha_{\mathrm{edge}}
\prod_{a \in \mathrm{Anc}(u)}
\frac{R_a}{|F(a)|},
\end{equation}
where $R_a$ is the number of rejections already made in the ancestor family
$F(a)$. Intuitively, evidence is allowed to propagate downward only after
earlier parts of the tree have already shown signal. Root families start at
$\alpha_{\mathrm{edge}}$. Within each family, ordinary Benjamini--Hochberg is
applied at level $\alpha_u$. Let
$S^{\mathrm{edge}}(c) \in \{0,1\}$ denote the resulting significance indicator
for child $c$.
